\section{Radical and Socle Series}

\begin{proposition}
    Let $A$ be an Artinian ring, and let $U$ be a finitely generated left $A$-module.

    \begin{enumerate}
        \item For a submodule $V\leq U$, the following are equivalent:
              \begin{enumerate}
                  \item $V=\operatorname{Rad}(U)$;
                  \item $V$ is the smallest submodule of $U$ such that $U/V$ is semisimple;
                  \item $V=J(A)U$.
              \end{enumerate}

        \item For a submodule $V\leq U$, the following are equivalent:
              \begin{enumerate}
                  \item $V=\operatorname{Soc}(U)$;
                  \item $V$ is the largest semisimple submodule of $U$;
                  \item
                        \[
                            V=\{\,u\in U\mid J(A)u=0\,\}.
                        \]
              \end{enumerate}
    \end{enumerate}
\end{proposition}

\begin{proof}
    We first record two standard facts.

    \medskip

    \noindent
    \textbf{Fact 1.}
    Since $A$ is Artinian, the quotient ring $A/J(A)$ is semisimple. Hence every
    $A/J(A)$-module is semisimple.

    \medskip

    \noindent
    \textbf{Fact 2.}
    If $S$ is a simple left $A$-module, then
    \[
        J(A)S=0.
    \]
    Indeed, by the characterization
    \[
        J(A)=\bigcap \{\,\operatorname{Ann}_A(T)\mid T \text{ simple left }A\text{-module}\,\},
    \]
    we have $J(A)\subseteq \operatorname{Ann}_A(S)$.

    \bigskip

    \noindent
    \textbf{(1)} We prove the equivalence of \textup{(a)}, \textup{(b)}, \textup{(c)}.

    \medskip

    \noindent
    \textit{\textup{(c)} $\Rightarrow$ \textup{(b)}.}
    Set
    \[
        W:=J(A)U.
    \]
    Then
    \[
        J(A)(U/W)=0,
    \]
    so $U/W$ is naturally an $A/J(A)$-module. By Fact~1, $U/W$ is semisimple.

    Now let $V'\leq U$ be any submodule such that $U/V'$ is semisimple. Since every
    simple module is annihilated by $J(A)$, every semisimple module is annihilated by
    $J(A)$ as well. Hence
    \[
        J(A)(U/V')=0.
    \]
    Equivalently,
    \[
        J(A)U\subseteq V'.
    \]
    Thus $W=J(A)U$ is the smallest submodule of $U$ with semisimple quotient.

    \medskip

    \noindent
    \textit{\textup{(b)} $\Rightarrow$ \textup{(c)}.}
    Assume that $V$ is the smallest submodule of $U$ such that $U/V$ is semisimple.

    Since $U/J(A)U$ is semisimple by the previous paragraph, the minimality of $V$ gives
    \[
        V\subseteq J(A)U.
    \]
    On the other hand, $U/V$ is semisimple, so it is annihilated by $J(A)$. Hence
    \[
        J(A)(U/V)=0,
    \]
    which means
    \[
        J(A)U\subseteq V.
    \]
    Therefore
    \[
        V=J(A)U.
    \]

    \medskip

    \noindent
    \textit{\textup{(c)} $\Leftrightarrow$ \textup{(a)}.}
    Let $W=J(A)U$. We show that $W=\operatorname{Rad}(U)$.

    First, let $M$ be a maximal submodule of $U$. Then $U/M$ is simple, hence semisimple.
    By the minimality proved above for $J(A)U$, we must have
    \[
        J(A)U\subseteq M.
    \]
    Since this is true for every maximal submodule $M$ of $U$, it follows that
    \[
        J(A)U\subseteq \operatorname{Rad}(U).
    \]

    For the reverse inclusion, suppose there exists
    \[
        u\in \operatorname{Rad}(U)\setminus J(A)U.
    \]
    Let
    \[
        \overline{U}:=U/J(A)U,\qquad \overline{u}:=u+J(A)U\neq 0.
    \]
    We already know that $\overline{U}$ is semisimple, so we may write
    \[
        \overline{U}=\bigoplus_{i\in I} S_i
    \]
    with each $S_i$ simple. Since $\overline{u}\neq 0$, some component of $\overline{u}$ in
    this direct sum is nonzero; choose $i_0\in I$ such that the $S_{i_0}$-component of
    $\overline{u}$ is nonzero.

    Set
    \[
        \overline{M}:=\bigoplus_{i\neq i_0} S_i.
    \]
    Then $\overline{M}$ is a maximal submodule of $\overline{U}$ and
    \[
        \overline{u}\notin \overline{M}.
    \]
    Let $M$ be the inverse image of $\overline{M}$ under the quotient map
    \[
        \pi:U\twoheadrightarrow \overline{U}.
    \]
    Then $M$ is a maximal submodule of $U$, because
    \[
        U/M\cong \overline{U}/\overline{M}\cong S_{i_0}
    \]
    is simple. But $u\notin M$, contradicting $u\in \operatorname{Rad}(U)$, since
    $\operatorname{Rad}(U)$ is contained in every maximal submodule of $U$.

    Hence
    \[
        \operatorname{Rad}(U)\subseteq J(A)U.
    \]
    Together with the opposite inclusion, we get
    \[
        \operatorname{Rad}(U)=J(A)U.
    \]

    Combining the implications above, we conclude that \textup{(a)}, \textup{(b)}, and
    \textup{(c)} are equivalent.

    \bigskip

    \noindent
    \textbf{(2)} Recall that
    \[
        \operatorname{Soc}(U)=\sum \{\,S\leq U\mid S \text{ is a simple submodule}\,\}.
    \]

    Let
    \[
        T:=\{\,u\in U\mid J(A)u=0\,\}.
    \]
    We first show that $T$ is a submodule of $U$. If $u,v\in T$ and $a\in A$, then
    \[
        J(A)(u+v)=J(A)u+J(A)v=0,
    \]
    and for every $j\in J(A)$,
    \[
        j(au)=(ja)u=0,
    \]
    because $J(A)$ is a two-sided ideal and $ja\in J(A)$. Thus $au\in T$, so $T\leq U$.

    \medskip

    \noindent
    \textit{\textup{(c)} $\Rightarrow$ \textup{(b)}.}
    Assume $V=T$.

    Since $J(A)T=0$, the module $T$ is naturally an $A/J(A)$-module. By Fact~1, $T$ is
    semisimple.

    Now let $W\leq U$ be any semisimple submodule. Then $W$ is a sum of simple
    submodules, and by Fact~2 each simple submodule is annihilated by $J(A)$. Hence
    \[
        J(A)W=0,
    \]
    so
    \[
        W\subseteq T.
    \]
    Therefore $T$ is the largest semisimple submodule of $U$.

    \medskip

    \noindent
    \textit{\textup{(b)} $\Rightarrow$ \textup{(c)}.}
    Assume that $V$ is the largest semisimple submodule of $U$.

    The submodule
    \[
        T=\{\,u\in U\mid J(A)u=0\,\}
    \]
    is semisimple by the previous paragraph, so by maximality of $V$ we get
    \[
        T\subseteq V.
    \]
    On the other hand, $V$ is semisimple, hence $J(A)V=0$, so
    \[
        V\subseteq T.
    \]
    Thus
    \[
        V=T.
    \]

    \medskip

    \noindent
    \textit{\textup{(a)} $\Leftrightarrow$ \textup{(c)}.}
    Let $T=\{u\in U\mid J(A)u=0\}$ as above.

    Every simple submodule $S\leq U$ satisfies $J(A)S=0$ by Fact~2, so
    \[
        S\subseteq T.
    \]
    Taking the sum over all simple submodules of $U$, we obtain
    \[
        \operatorname{Soc}(U)\subseteq T.
    \]

    Conversely, $T$ is semisimple, so it is a sum of simple submodules of $U$. Hence
    \[
        T\subseteq \operatorname{Soc}(U).
    \]
    Therefore
    \[
        \operatorname{Soc}(U)=T.
    \]

    So \textup{(a)}, \textup{(b)}, and \textup{(c)} are equivalent.
\end{proof}

Let $A$ be an Artinian ring, and let $U$ be a finitely generated left $A$-module.
Recall the basic identities
\[
    \operatorname{Rad}(U)=J(A)U,
    \qquad
    \operatorname{Soc}(U)=\{u\in U\mid J(A)u=0\}.
\]

\begin{definition}
    Define inductively
    \[
        \operatorname{Rad}^0(U):=U,\qquad \operatorname{Rad}^n(U):=\operatorname{Rad}(\operatorname{Rad}^{n-1}(U))\quad (n\ge 1),
    \]
    and
    \[
        \operatorname{Soc}^0(U):=0,\qquad
        \operatorname{Soc}^n(U)/\operatorname{Soc}^{n-1}(U):=\operatorname{Soc}\bigl(U/\operatorname{Soc}^{n-1}(U)\bigr)\quad (n\ge 1).
    \]
    The descending chain
    \[
        U=\operatorname{Rad}^0(U)\supseteq \operatorname{Rad}(U)\supseteq \operatorname{Rad}^2(U)\supseteq \cdots
    \]
    is called the \emph{radical series} of $U$, and the ascending chain
    \[
        0=\operatorname{Soc}^0(U)\subseteq \operatorname{Soc}(U)\subseteq \operatorname{Soc}^2(U)\subseteq \cdots
    \]
    is called the \emph{socle series} of $U$.
\end{definition}

\begin{proposition}
    For every integer $n\ge 1$,\label{prop:rad-soc-formulas}
    \[
        \operatorname{Rad}^n(U)=J(A)^nU,
    \]
    and
    \[
        \operatorname{Soc}^n(U)=\{u\in U\mid J(A)^n u=0\}.
    \]
\end{proposition}

\begin{proof}
    We prove both formulas by induction on $n$.

    \medskip

    \noindent
    \textbf{Radical formula.}
    For $n=1$, this is exactly the identity
    \[
        \operatorname{Rad}(U)=J(A)U.
    \]
    Assume now that $\operatorname{Rad}^{n-1}(U)=J(A)^{n-1}U$. Since $\operatorname{Rad}^{n-1}(U)$ is a finitely generated
    $A$-module, we may apply the same identity to it:
    \[
        \operatorname{Rad}^n(U)
        =\operatorname{Rad}(\operatorname{Rad}^{n-1}(U))
        =J(A)\operatorname{Rad}^{n-1}(U)
        =J(A)\bigl(J(A)^{n-1}U\bigr)
        =J(A)^nU.
    \]

    \medskip

    \noindent
    \textbf{Socle formula.}
    For $n=1$, this is exactly the identity
    \[
        \operatorname{Soc}(U)=\{u\in U\mid J(A)u=0\}.
    \]
    Assume now that
    \[
        \operatorname{Soc}^{n-1}(U)=\{u\in U\mid J(A)^{n-1}u=0\}.
    \]
    By definition,
    \[
        \operatorname{Soc}^n(U)/\operatorname{Soc}^{n-1}(U)
        =
        \operatorname{Soc}\bigl(U/\operatorname{Soc}^{n-1}(U)\bigr).
    \]
    Applying the case $n=1$ to the module $U/\operatorname{Soc}^{n-1}(U)$, we get
    \[
        \operatorname{Soc}\bigl(U/\operatorname{Soc}^{n-1}(U)\bigr)
        =
        \left\{
        u+\operatorname{Soc}^{n-1}(U)\in U/\operatorname{Soc}^{n-1}(U)
        \;\middle|\;
        J(A)u\subseteq \operatorname{Soc}^{n-1}(U)
        \right\}.
    \]
    By the induction hypothesis, the condition
    \[
        J(A)u\subseteq \operatorname{Soc}^{n-1}(U)
    \]
    is equivalent to
    \[
        J(A)^{n-1}(J(A)u)=0,
    \]
    that is,
    \[
        J(A)^n u=0.
    \]
    Hence the preimage of $\operatorname{Soc}(U/\operatorname{Soc}^{n-1}(U))$ in $U$ is precisely
    \[
        \{u\in U\mid J(A)^n u=0\}.
    \]
    Therefore
    \[
        \operatorname{Soc}^n(U)=\{u\in U\mid J(A)^n u=0\}.
    \]
    This completes the induction.
\end{proof}

\begin{remark}
    For $n\ge 1$, the quotient
    \[
        \operatorname{Rad}^{n-1}(U)/\operatorname{Rad}^n(U)
    \]
    is called the $n$-th \emph{radical layer} of $U$, and
    \[
        \operatorname{Soc}^n(U)/\operatorname{Soc}^{n-1}(U)
    \]
    is called the $n$-th \emph{socle layer} of $U$.
\end{remark}

\begin{remark}
    The $n$-th radical layer is the top of the previous radical term:
    \[
        \operatorname{Rad}^{n-1}(U)/\operatorname{Rad}^{n}(U)
        =
        \operatorname{Top}\bigl(\operatorname{Rad}^{n-1}(U)\bigr).
    \]
\end{remark}

\begin{lemma}\label{lem:rad-soc-direct-sum}
    For every $n\ge 1$ and any finitely generated $A$-modules $U,V$,
    \[
        \operatorname{Rad}^n(U\oplus V)=\operatorname{Rad}^n(U)\oplus \operatorname{Rad}^n(V),
    \]
    and
    \[
        \operatorname{Soc}^n(U\oplus V)=\operatorname{Soc}^n(U)\oplus \operatorname{Soc}^n(V).
    \]
\end{lemma}

\begin{proof}
    Using Proposition~\ref{prop:rad-soc-formulas},
    \[
        \operatorname{Rad}^n(U\oplus V)
        =
        J(A)^n(U\oplus V)
        =
        J(A)^nU\oplus J(A)^nV
        =
        \operatorname{Rad}^n(U)\oplus \operatorname{Rad}^n(V).
    \]
    Similarly,
    \[
        \operatorname{Soc}^n(U\oplus V)
        =
        \{(u,v)\in U\oplus V\mid J(A)^n(u,v)=0\}.
    \]
    But
    \[
        J(A)^n(u,v)=0
        \iff
        J(A)^n u=0 \text{ and } J(A)^n v=0.
    \]
    Hence
    \[
        \operatorname{Soc}^n(U\oplus V)
        =
        \{u\in U\mid J(A)^n u=0\}
        \oplus
        \{v\in V\mid J(A)^n v=0\}
        =
        \operatorname{Soc}^n(U)\oplus \operatorname{Soc}^n(V).
    \]
\end{proof}

\begin{theorem}\label{thm:rad-soc-fastest}
    Let $A$ be an Artinian ring, and let $U$ be a finitely generated left $A$-module.

    \begin{enumerate}
        \item The radical series of $U$ is the fastest descending chain
              \[
                  U=U_0\supseteq U_1\supseteq U_2\supseteq \cdots
              \]
              such that every factor $U_{r-1}/U_r$ is semisimple. More precisely, if
              \[
                  U=U_0\supseteq U_1\supseteq U_2\supseteq \cdots
              \]
              is any descending chain with each $U_{r-1}/U_r$ semisimple, then for every $r\ge 0$,
              \[
                  \operatorname{Rad}^r(U)\subseteq U_r.
              \]

        \item The socle series of $U$ is the fastest ascending chain
              \[
                  0=V_0\subseteq V_1\subseteq V_2\subseteq \cdots
              \]
              such that every factor $V_r/V_{r-1}$ is semisimple. More precisely, if
              \[
                  0=V_0\subseteq V_1\subseteq V_2\subseteq \cdots
              \]
              is any ascending chain with each $V_r/V_{r-1}$ semisimple, then for every $r\ge 0$,
              \[
                  V_r\subseteq \operatorname{Soc}^r(U).
              \]

        \item Both series terminate. If $m$ and $n$ are the minimal integers such that
              \[
                  \operatorname{Rad}^m(U)=0
                  \qquad\text{and}\qquad
                  \operatorname{Soc}^n(U)=U,
              \]
              then
              \[
                  m=n.
              \]
    \end{enumerate}
\end{theorem}

\begin{proof}
    \textbf{(1)} Let
    \[
        U=U_0\supseteq U_1\supseteq U_2\supseteq \cdots
    \]
    be a descending chain such that each quotient $U_{r-1}/U_r$ is semisimple.
    We prove by induction on $r$ that
    \[
        \operatorname{Rad}^r(U)\subseteq U_r.
    \]

    For $r=0$, this is clear, since $\operatorname{Rad}^0(U)=U=U_0$.

    Assume $r\ge 1$ and that
    \[
        \operatorname{Rad}^{r-1}(U)\subseteq U_{r-1}.
    \]
    Consider the quotient
    \[
        \operatorname{Rad}^{r-1}(U)\big/\bigl(\operatorname{Rad}^{r-1}(U)\cap U_r\bigr).
    \]
    By the second isomorphism theorem,
    \[
        \operatorname{Rad}^{r-1}(U)\big/\bigl(\operatorname{Rad}^{r-1}(U)\cap U_r\bigr)
        \cong
        \bigl(\operatorname{Rad}^{r-1}(U)+U_r\bigr)/U_r.
    \]
    Since $\operatorname{Rad}^{r-1}(U)\subseteq U_{r-1}$, we have
    \[
        \bigl(\operatorname{Rad}^{r-1}(U)+U_r\bigr)/U_r
        \subseteq
        U_{r-1}/U_r.
    \]
    But $U_{r-1}/U_r$ is semisimple, and any submodule of a semisimple module is semisimple.
    Hence
    \[
        \operatorname{Rad}^{r-1}(U)\big/\bigl(\operatorname{Rad}^{r-1}(U)\cap U_r\bigr)
    \]
    is semisimple.

    By definition of the radical as the smallest submodule with semisimple quotient, it follows that
    \[
        \operatorname{Rad}\bigl(\operatorname{Rad}^{r-1}(U)\bigr)
        \subseteq
        \operatorname{Rad}^{r-1}(U)\cap U_r.
    \]
    That is,
    \[
        \operatorname{Rad}^r(U)\subseteq U_r.
    \]
    This proves \textup{(1)}.

    \bigskip

    \textbf{(2)} Let
    \[
        0=V_0\subseteq V_1\subseteq V_2\subseteq \cdots
    \]
    be an ascending chain such that each quotient $V_r/V_{r-1}$ is semisimple.
    We prove by induction on $r$ that
    \[
        V_r\subseteq \operatorname{Soc}^r(U).
    \]

    For $r=0$, this is clear.

    Assume $r\ge 1$ and that
    \[
        V_{r-1}\subseteq \operatorname{Soc}^{r-1}(U).
    \]
    Consider the image of $V_r$ in the quotient $U/\operatorname{Soc}^{r-1}(U)$:
    \[
        \bigl(V_r+\operatorname{Soc}^{r-1}(U)\bigr)\big/\operatorname{Soc}^{r-1}(U).
    \]
    By the second isomorphism theorem,
    \[
        \bigl(V_r+\operatorname{Soc}^{r-1}(U)\bigr)\big/\operatorname{Soc}^{r-1}(U)
        \cong
        V_r/\bigl(V_r\cap \operatorname{Soc}^{r-1}(U)\bigr).
    \]
    Since $V_{r-1}\subseteq \operatorname{Soc}^{r-1}(U)$, we have
    \[
        V_{r-1}\subseteq V_r\cap \operatorname{Soc}^{r-1}(U),
    \]
    so
    \[
        V_r/\bigl(V_r\cap \operatorname{Soc}^{r-1}(U)\bigr)
    \]
    is a quotient of $V_r/V_{r-1}$. Because $V_r/V_{r-1}$ is semisimple, so is every quotient of it.
    Hence
    \[
        \bigl(V_r+\operatorname{Soc}^{r-1}(U)\bigr)\big/\operatorname{Soc}^{r-1}(U)
    \]
    is a semisimple submodule of
    \[
        U/\operatorname{Soc}^{r-1}(U).
    \]
    By the maximality of the socle,
    \[
        \bigl(V_r+\operatorname{Soc}^{r-1}(U)\bigr)\big/\operatorname{Soc}^{r-1}(U)
        \subseteq
        \operatorname{Soc}\bigl(U/\operatorname{Soc}^{r-1}(U)\bigr)
        =
        \operatorname{Soc}^r(U)/\operatorname{Soc}^{r-1}(U).
    \]
    Therefore
    \[
        V_r\subseteq \operatorname{Soc}^r(U).
    \]
    This proves \textup{(2)}.

    \bigskip

    \textbf{(3)} Since $A$ is Artinian, $J(A)$ is nilpotent. Choose $N\ge 1$ such that
    \[
        J(A)^N=0.
    \]
    Then by Proposition~\ref{prop:rad-soc-formulas},
    \[
        \operatorname{Rad}^N(U)=J(A)^N U=0,
    \]
    and
    \[
        \operatorname{Soc}^N(U)=\{u\in U\mid J(A)^N u=0\}=U.
    \]
    So both series terminate.

    Now let $m$ be minimal such that $\operatorname{Rad}^m(U)=0$. By Proposition~\ref{prop:rad-soc-formulas},
    \[
        J(A)^mU=0.
    \]
    Hence every $u\in U$ satisfies $J(A)^m u=0$, so again by
    Proposition~\ref{prop:rad-soc-formulas},
    \[
        \operatorname{Soc}^m(U)=U.
    \]
    Thus $n\le m$.

    Conversely, let $n$ be minimal such that $\operatorname{Soc}^n(U)=U$. Then for every $u\in U$,
    \[
        J(A)^n u=0.
    \]
    So
    \[
        J(A)^nU=0,
    \]
    hence
    \[
        \operatorname{Rad}^n(U)=J(A)^nU=0.
    \]
    Therefore $m\le n$.

    Combining the two inequalities, we obtain
    \[
        m=n.
    \]
\end{proof}

\begin{definition}
    The common integer in Theorem~\ref{thm:rad-soc-fastest}(3) is called the
    \emph{Loewy length} of $U$.
\end{definition}

